\begin{center}
Problema J 

Álbum da Copa

{\tiny Por OBI BR Brasil}

Timelimit: $1$	
\end{center}			
			
Em ano de Copa do Mundo de Futebol, o álbum de figurinhas oficial é sempre um grande sucesso entre crianças e também entre adultos. Para quem não conhece, o álbum contém espaços numerados de 1 a N para colar as figurinhas; cada figurinha, também numerada de 1 a N, é uma pequena foto de um jogador de uma das seleções que jogará a Copa do Mundo. O objetivo é colar todas as figurinhas nos respectivos espaços no álbum, de modo a completar o álbum (ou seja, não deixar nenhum espaço sem a correspondente figurinha).

As figurinhas são vendidas em envelopes fechados, de forma que o comprador não sabe quais figurinhas está comprando, e pode ocorrer de comprar uma figurinha que ele já tenha colado no álbum.

Para ajudar os usuários, a empresa responsável pela venda do álbum e das figurinhas quer criar um aplicativo que permita gerenciar facilmente as figurinhas que faltam para completar o álbum e está solicitando a sua ajuda.

Dados o número total de espaços e figurinhas do álbum, e uma lista das figurinhas já compradas (que pode conter figurinhas repetidas), sua tarefa é determinar quantas figurinhas faltam para completar o álbum.
 
\vspace{0.3cm}
\textbf{Entrada}
A primeira linha contém um inteiro $N (1 \leq N \leq 100)$ indicando o número total de figurinhas e espaços no álbum. A segunda linha contém um inteiro $M (1 \leq M \leq 300)$ indicando o número de figurinhas já compradas. Cada uma das $M$ linhas seguintes contém um número inteiro $X (1 \leq X \leq N)$ indicando uma figurinha já comprada.

\vspace{0.3cm}
\textbf{Saída}
Seu programa deve produzir uma única linha contendo um inteiro representando o número de figurinhas que falta para completar o álbum.
\begin{table}[h]
    \centering
    \begin{tabular}{|p{7cm}|p{7cm}|}
        \hline
         \begin{center}
            \vspace{-0.3cm}
             \textbf{Exemplos de Entrada}
            \vspace{-0.3cm}
         \end{center} 
         &  
         \begin{center}
            \vspace{-0.3cm}
             \textbf{Exemplos de Saída}
            \vspace{-0.3cm}
         \end{center} \\
        \hline
         \texttt{10 3} & \texttt{7} \\
         \texttt{5 8 3} & \\
        \hline
         \texttt{5 6} & \texttt{3} \\
         \texttt{3 3 2 3 3 3} & \\
         \hline         
          \texttt{3 4} & \texttt{0} \\
         \texttt{2 1 3 3} & \\
         \hline             
    \end{tabular}
    \caption{OBI - Olimpíada Brasileira de Informática 2018 - Fase 1}
    \label{tab:inputb}
    
\end{table}

\newpage

\begin{center}
\noindent
\lstinputlisting{abobora_22/J/J2.cpp}
\end{center}